\documentclass{article}
\usepackage{amsmath,amssymb,amsthm}
\usepackage{algorithm}
\usepackage[noend]{algpseudocode}
\usepackage{enumitem}
\usepackage{hyperref}
\usepackage{geometry}
\geometry{margin=1in}
\newtheorem{theorem}{Theorem}
\newtheorem{lemma}{Lemma}
\newtheorem{assumption}{Assumption}
\newcommand{\norm}[1]{\left\lVert#1\right\rVert}
\newcommand{\R}{\mathbb{R}}

\begin{document}

\section*{Algorithm Name}
\textbf{Trust‑Region Method with Approximate Subproblem Solution (TR‑A)}  

The method iteratively builds a quadratic model of a smooth non‑convex objective
\(f:\R^{n}\rightarrow\R\) around the current iterate \(x_{k}\) and restricts the
step to a trust region of radius \(\Delta_{k}>0\).  The step \(p_{k}\) is required
only to satisfy a \emph{sufficient‑decrease condition} rather than to be the exact
solution of the subproblem.

\section*{Mathematical Formulation}
At iteration \(k\) let
\[
g_{k}= \nabla f(x_{k}),\qquad 
B_{k}= \nabla^{2}f(x_{k})\;( \text{or a symmetric approximation} ).
\]

The quadratic model is
\[
m_{k}(p)=f(x_{k})+g_{k}^{\top}p+\tfrac12 p^{\top}B_{k}p .
\]

The trust‑region subproblem is
\[
\min_{p\in\R^{n}}\; m_{k}(p)\quad\text{s.t.}\quad \norm{p}\le \Delta_{k}. \tag{TR}
\]

Instead of solving (TR) exactly we compute an \emph{approximate} step
\(p_{k}\) that satisfies the \textbf{Cauchy decrease condition}
\[
m_{k}(0)-m_{k}(p_{k})\;\ge\;\eta_{1}\bigl(m_{k}(0)-m_{k}(p_{k}^{C})\bigr),\qquad
0<\eta_{1}<1,
\tag{1}
\]
where \(p_{k}^{C}\) is the Cauchy point (the minimizer of the model along the
steepest‑descent direction within the trust region).

Define the actual reduction ratio
\[
\rho_{k}= \frac{f(x_{k})-f(x_{k}+p_{k})}{m_{k}(0)-m_{k}(p_{k})}.
\tag{2}
\]

The trust‑region radius is updated according to two thresholds
\(0<\eta_{1}<\eta_{2}<1\):
\[
\Delta_{k+1}=
\begin{cases}
\gamma_{1}\Delta_{k}, & \rho_{k}<\eta_{1},\\[4pt]
\Delta_{k},           & \eta_{1}\le\rho_{k}<\eta_{2},\\[4pt]
\gamma_{2}\Delta_{k}, & \rho_{k}\ge\eta_{2},
\end{cases}
\tag{3}
\]
with \(0<\gamma_{1}<1<\gamma_{2}\).

The new iterate is
\[
x_{k+1}= 
\begin{cases}
x_{k}+p_{k}, & \rho_{k}\ge \eta_{1},\\
x_{k},       & \rho_{k}<\eta_{1}.
\end{cases}
\tag{4}
\]

\section*{Pseudocode}
\begin{algorithm}[H]
\caption{Trust‑Region Method with Approximate Subproblem Solution (TR‑A)}
\begin{algorithmic}[1]
\Require Initial point \(x_{0}\), initial radius \(\Delta_{0}>0\),
          parameters \(\eta_{1},\eta_{2},\gamma_{1},\gamma_{2}\) as in (3),
          tolerance \(\varepsilon>0\).
\State \(k\gets 0\)
\While{\(\norm{\nabla f(x_{k})}> \varepsilon\)}
    \State Compute \(g_{k}\gets\nabla f(x_{k})\) and a symmetric \(B_{k}\) (e.g. Hessian).
    \State Obtain an approximate step \(p_{k}\) satisfying (1).   \Comment{e.g. truncated CG}
    \State Compute the ratio \(\rho_{k}\) by (2).
    \If{\(\rho_{k}\ge \eta_{1}\)} 
        \State \(x_{k+1}\gets x_{k}+p_{k}\)          \Comment{accept step}
    \Else
        \State \(x_{k+1}\gets x_{k}\)                \Comment{reject step}
    \EndIf
    \State Update \(\Delta_{k+1}\) using (3).
    \State \(k\gets k+1\)
\EndWhile
\State \Return \(x_{k}\)
\end{algorithmic}
\end{algorithm}

\section*{Dry Run Example}
Consider the one‑dimensional quadratic
\[
f(w)=(w-3)^{2},\qquad w\in\R .
\]
We set
\[
w_{0}=0,\quad \Delta_{0}=1,\quad 
\eta_{1}=0.1,\;\eta_{2}=0.75,\;
\gamma_{1}=0.5,\;\gamma_{2}=2.0 .
\]

Because the problem is 1‑D, the Cauchy point coincides with the exact solution
of (TR).  The gradient and Hessian are
\[
g_{k}=2(w_{k}-3),\qquad B_{k}=2.
\]

The exact minimizer of the model under the trust‑region bound is
\[
p_{k}^{*}= -\frac{g_{k}}{B_{k}}=-(w_{k}-3).
\]
If \(|p_{k}^{*}|\le\Delta_{k}\) we take \(p_{k}=p_{k}^{*}\); otherwise we truncate
to \(p_{k}= -\operatorname{sign}(g_{k})\Delta_{k}\).

\medskip
\textbf{Iteration 0}
\[
\begin{aligned}
w_{0}&=0, &
g_{0}&=2(0-3)=-6, &
p_{0}^{*}&=-(0-3)=3. 
\end{aligned}
\]
Since \(|p_{0}^{*}|=3>\Delta_{0}=1\), we truncate:
\(p_{0}=-\operatorname{sign}(-6)\cdot1= -1\).

Model reduction:
\(m_{0}(0)-m_{0}(p_{0}) = -g_{0}p_{0}-\tfrac12 B_{0}p_{0}^{2}
=6\cdot1-\tfrac12\cdot2\cdot1^{2}=5.\)

Actual reduction:
\(f(w_{0})-f(w_{0}+p_{0}) = (0-3)^{2}-( -1-3)^{2}=9-16=-7\) (increase!).
Thus
\(\rho_{0}=(-7)/5=-1.4<\eta_{1}\).
The step is rejected, \(w_{1}=w_{0}=0\), and the radius is shrunk:
\(\Delta_{1}= \gamma_{1}\Delta_{0}=0.5\).

\medskip
\textbf{Iteration 1}
\[
\Delta_{1}=0.5,\quad p_{1}^{*}=3\;(>0.5)\Rightarrow p_{1}=-0.5.
\]
Model reduction:
\(m_{1}(0)-m_{1}(p_{1}) = 6\cdot0.5-\tfrac12\cdot2\cdot0.5^{2}=3-0.25=2.75.\)

Actual reduction:
\(f(0)-f(-0.5)=9-( -0.5-3)^{2}=9-12.25=-3.25.\)

\(\rho_{1}=(-3.25)/2.75\approx-1.18<\eta_{1}\) ⇒ reject,
\(w_{2}=0\), \(\Delta_{2}=0.25\).

\medskip
\textbf{Iteration 2}
\[
\Delta_{2}=0.25,\; p_{2}=-0.25.
\]
Model reduction:
\(m_{2}(0)-m_{2}(p_{2}) = 6\cdot0.25-\tfrac12\cdot2\cdot0.25^{2}=1.5-0.0625=1.4375.\)

Actual reduction:
\(f(0)-f(-0.25)=9-( -0.25-3)^{2}=9-10.5625=-1.5625.\)

\(\rho_{2}=(-1.5625)/1.4375\approx-1.09<\eta_{1}\) ⇒ reject,
\(w_{3}=0\), \(\Delta_{3}=0.125\).

The example illustrates that, when the initial radius is too small,
the algorithm keeps rejecting steps and shrinks the trust region.
If we start with a larger \(\Delta_{0}\) (e.g. \(\Delta_{0}=5\)), the first
truncated step would be \(p_{0}= -3\) and would be accepted, driving the
iterates directly to the minimizer \(w^{*}=3\).

\newpage
\section*{Assumptions}
\begin{assumption}[Smoothness]\label{ass:smooth}
\(f\) is twice continuously differentiable and its gradient is \(L\)-Lipschitz:
\[
\|\nabla f(x)-\nabla f(y)\|\le L\|x-y\|,\qquad\forall x,y\in\R^{n}.
\]
\end{assumption}
\begin{assumption}[Bounded Hessian]\label{ass:hessian}
There exists \(M>0\) such that
\[
\|B_{k}\|\le M,\qquad\forall k,
\]
where \(\|\cdot\|\) denotes the spectral norm.
\end{assumption}
\begin{assumption}[Level‑set boundedness]\label{ass:level}
For the initial point \(x_{0}\) the set
\[
\mathcal{L}:=\{x\in\R^{n}\mid f(x)\le f(x_{0})\}
\]
is compact.
\end{assumption}
\begin{assumption}[Approximate solution]\label{ass:approx}
The step \(p_{k}\) returned by the subproblem solver satisfies the
Cauchy decrease condition (1) with a fixed constant \(\eta_{1}\in(0,1)\).
\end{assumption}

\section*{Main Convergence Theorem}
\begin{theorem}[Global convergence and worst‑case complexity]\label{thm:main}
Suppose Assumptions \ref{ass:smooth}--\ref{ass:approx} hold and the
parameters satisfy
\[
0<\eta_{1}<\eta_{2}<1,\qquad 0<\gamma_{1}<1<\gamma_{2}.
\]
Then the sequence \(\{x_{k}\}\) generated by TR‑A satisfies
\[
\liminf_{k\to\infty}\|\nabla f(x_{k})\|=0.
\]
Moreover, for any tolerance \(\varepsilon>0\) the number of iterations
required to obtain a point with \(\|\nabla f(x_{k})\|\le\varepsilon\) is bounded by
\[
\mathcal{O}\!\left(\frac{L\, (f(x_{0})-f^{\star})}{\varepsilon^{2}}\right),
\tag{5}
\]
where \(f^{\star}=\inf_{x}f(x)\).
\end{theorem}

\section*{Theoretical Properties}
\begin{itemize}[leftmargin=*]
\item \textbf{Stability.}  Because the radius is reduced whenever the model
prediction is poor (\(\rho_{k}<\eta_{1}\)), the algorithm never makes steps that
increase the objective by more than a factor proportional to the model error.
\item \textbf{Hyper‑parameter sensitivity.}
  The constants \(\eta_{1},\eta_{2},\gamma_{1},\gamma_{2}\) only affect the
  constants hidden in the \(\mathcal{O}(\cdot)\) bound; the asymptotic rate
  \(\varepsilon^{-2}\) is independent of their precise values as long as they
  respect the ordering in Theorem~\ref{thm:main}.
\item \textbf{Comparison to baselines.}
  For smooth non‑convex problems, SGD with diminishing stepsizes also attains an
  \(\varepsilon^{-2}\) iteration complexity in terms of expected gradient norm.
  The trust‑region method, however, provides deterministic guarantees and
  automatically adapts the step size via \(\Delta_{k}\), which can be beneficial
  when the Lipschitz constant \(L\) is unknown.
\end{itemize}

\section*{Discussion}
The result mirrors classical trust‑region theory (e.g. \cite{conn1999}).  The
novelty of the present analysis lies in the explicit use of the Cauchy
decrease condition (1) which is easy to verify for practical inner solvers
such as truncated conjugate‑gradient.  The bound (5) shows that the algorithm
behaves like a first‑order method in the worst case, but it can exploit
second‑order information when the model is accurate, often leading to
substantially fewer iterations in practice.

\newpage
\section*{Preliminaries (Definitions and Lemmas)}
\begin{lemma}[Model error bound]\label{lem:modelerror}
Under Assumption \ref{ass:smooth} and \ref{ass:hessian},
for any \(p\) with \(\|p\|\le\Delta_{k}\) we have
\[
\bigl|f(x_{k}+p)-m_{k}(p)\bigr|
\le \frac{L}{6}\|p\|^{3}+\frac{M}{2}\|p\|^{2}.
\tag{L1}
\]
\end{lemma}
\begin{proof}
By Taylor’s theorem with remainder,
\[
f(x_{k}+p)=f(x_{k})+g_{k}^{\top}p+\tfrac12 p^{\top}B_{k}p+r_{k}(p),
\]
where \(\|r_{k}(p)\|\le \frac{L}{6}\|p\|^{3}\) because the gradient is \(L\)‑Lipschitz.
Since \(\|B_{k}\|\le M\), the quadratic term error is bounded by \(\frac{M}{2}\|p\|^{2}\).
Subtracting the model \(m_{k}(p)\) yields (L1). \qedhere
\end{proof}

\begin{lemma}[Cauchy decrease]\label{lem:cauchy}
Let \(p_{k}^{C}= -\frac{g_{k}^{\top}g_{k}}{g_{k}^{\top}B_{k}g_{k}}\,g_{k}\) if
\(\|g_{k}\|>0\) and \(\|p_{k}^{C}\|\le\Delta_{k}\); otherwise
\(p_{k}^{C}= -\frac{\Delta_{k}}{\|g_{k}\|}g_{k}\).  Then
\[
m_{k}(0)-m_{k}(p_{k}^{C})\;\ge\;
\frac{\|g_{k}\|^{2}}{2\,(M+\frac{\|g_{k}\|}{\Delta_{k}})}.
\tag{L2}
\]
\end{lemma}
\begin{proof}
The Cauchy point is the minimizer of the model along the direction \(-g_{k}\)
subject to the radius constraint.  Direct algebra (see e.g. \cite{conn1999})
gives the stated bound after using \(\|B_{k}\|\le M\). \qedhere
\end{proof}

\section*{Main Proof (Step‑by‑Step Derivation)}
We prove Theorem~\ref{thm:main}.  Throughout the proof we keep the
iteration index \(k\) arbitrary but fixed.

\medskip
\noindent\textbf{[Step 1]} (\emph{Relation between actual and predicted reduction}).  
From the definition of \(\rho_{k}\) (2) and Lemma~\ref{lem:modelerror},
\[
\begin{aligned}
1-\rho_{k}
&= \frac{m_{k}(0)-m_{k}(p_{k})-\bigl[f(x_{k})-f(x_{k}+p_{k})\bigr]}
        {m_{k}(0)-m_{k}(p_{k})}\\[4pt]
&\le \frac{|f(x_{k}+p_{k})-m_{k}(p_{k})|}
           {m_{k}(0)-m_{k}(p_{k})}
 \;\le\;
 \frac{\frac{L}{6}\|p_{k}\|^{3}+\frac{M}{2}\|p_{k}\|^{2}}
      {m_{k}(0)-m_{k}(p_{k})}.      \tag{S1}
\end{aligned}
\]

\noindent\textbf{[Step 2]} (\emph{Lower bound on the denominator}).  
Because the step satisfies the Cauchy condition (1),
\[
m_{k}(0)-m_{k}(p_{k})\;\ge\;
\eta_{1}\bigl(m_{k}(0)-m_{k}(p_{k}^{C})\bigr).
\tag{S2}
\]
Applying Lemma~\ref{lem:cauchy} to the right‑hand side yields
\[
m_{k}(0)-m_{k}(p_{k})\;\ge\;
\eta_{1}\frac{\|g_{k}\|^{2}}
               {2\bigl(M+\frac{\|g_{k}\|}{\Delta_{k}}\bigr)}.
\tag{S3}
\]

\noindent\textbf{[Step 3]} (\emph{Bounding \(\|p_{k}\|\)}).  
Since the step is feasible,
\(\|p_{k}\|\le\Delta_{k}\).  Insert this into (S1):
\[
1-\rho_{k}
\le
\frac{\frac{L}{6}\Delta_{k}^{3}+\frac{M}{2}\Delta_{k}^{2}}
     {m_{k}(0)-m_{k}(p_{k})}.
\tag{S4}
\]

\noindent\textbf{[Step 4]} (\emph{If the step is rejected}).  
When \(\rho_{k}<\eta_{1}\) the algorithm rejects the step and shrinks the
radius: \(\Delta_{k+1}=\gamma_{1}\Delta_{k}\).  Using (S4) we obtain
\[
\Delta_{k+1}
\le \gamma_{1}\Delta_{k}
\le \gamma_{1}\Bigl(\frac{L}{6\eta_{1}}+\frac{M}{2\eta_{1}}\Bigr)^{\!1/2}
      \frac{\|p_{k}\|^{1/2}}{\|g_{k}\|^{1/2}}.
\tag{S5}
\]
Thus the radius cannot shrink indefinitely without the gradient norm
decreasing.

\noindent\textbf{[Step 5]} (\emph{If the step is accepted}).  
When \(\rho_{k}\ge\eta_{1}\) we set \(x_{k+1}=x_{k}+p_{k}\) and obtain from the
definition of \(\rho_{k}\):
\[
f(x_{k})-f(x_{k+1})
= \rho_{k}\bigl(m_{k}(0)-m_{k}(p_{k})\bigr)
\ge \eta_{1}\bigl(m_{k}(0)-m_{k}(p_{k})\bigr).
\tag{S6}
\]
Insert the lower bound (S3) into (S6):
\[
f(x_{k})-f(x_{k+1})
\ge
\eta_{1}^{2}\,
\frac{\|g_{k}\|^{2}}
     {2\bigl(M+\frac{\|g_{k}\|}{\Delta_{k}}\bigr)}.
\tag{S7}
\]

\noindent\textbf{[Step 6]} (\emph{Summation over successful iterations}).  
Let \(\mathcal{S}\) denote the set of indices where the step is accepted.
Summing (S7) over \(k\in\mathcal{S}\) yields
\[
\sum_{k\in\mathcal{S}}
\frac{\|g_{k}\|^{2}}
     {M+\frac{\|g_{k}\|}{\Delta_{k}}}
\;\le\;
\frac{2}{\eta_{1}^{2}}\bigl(f(x_{0})-f^{\star}\bigr).
\tag{S8}
\]

\noindent\textbf{[Step 7]} (\emph{Bounding the denominator by a constant}).  
Since \(\Delta_{k}\le\Delta_{\max}\) for all \(k\) (the radius is never increased
beyond \(\gamma_{2}\Delta_{\max}\) and the level set is compact), we have
\[
M+\frac{\|g_{k}\|}{\Delta_{k}}
\le M+\frac{\|g_{k}\|}{\Delta_{\min}}
\le M+\frac{G}{\Delta_{\min}},
\]
where \(G\) is an upper bound on the gradient norm over the compact level set.
Thus
\[
\frac{1}{M+\frac{\|g_{k}\|}{\Delta_{k}}}
\ge \frac{1}{M+G/\Delta_{\min}}=:c>0.
\tag{S9}
\]

\noindent\textbf{[Step 8]} (\emph{Deriving the complexity bound}).  
Insert (S9) into (S8):
\[
c\sum_{k\in\mathcal{S}}\|g_{k}\|^{2}
\le
\frac{2}{\eta_{1}^{2}}\bigl(f(x_{0})-f^{\star}\bigr).
\]
Hence the number of successful iterations with \(\|g_{k}\|\ge\varepsilon\) is at most
\[
N_{\text{succ}}(\varepsilon)
\le
\frac{2\bigl(f(x_{0})-f^{\star}\bigr)}
     {c\,\eta_{1}^{2}\,\varepsilon^{2}}.
\tag{S10}
\]

\noindent\textbf{[Step 9]} (\emph{Relating successful to total iterations}).  
Between two successive successful steps the algorithm may perform at most
\(K:=\bigl\lceil\log_{\gamma_{1}}(\Delta_{\min}/\Delta_{\max})\bigr\rceil\)
rejections (because each rejection multiplies \(\Delta_{k}\) by \(\gamma_{1}<1\)
until it becomes small enough to satisfy \(\rho_{k}\ge\eta_{1}\)).  Therefore
the total number of iterations is bounded by a constant factor of
\(N_{\text{succ}}(\varepsilon)\).  Combining with (S10) yields the
complexity (5).

\noindent\textbf{[Step 10]} (\emph{Limit‑point result}).  
Since the sum of \(\|g_{k}\|^{2}\) over all successful indices is finite,
\(\|g_{k}\|\to0\) along a subsequence.  The radius updates guarantee that any
limit point of \(\{x_{k}\}\) satisfies the first‑order optimality condition,
which proves \(\displaystyle\liminf_{k\to\infty}\|\nabla f(x_{k})\|=0\).

\qed

\section*{Rate Derivation (With Constants)}
From (S10) we may write explicitly
\[
N(\varepsilon)
\;\le\;
\underbrace{\frac{2}{c\,\eta_{1}^{2}}}_{C_{1}}
\bigl(f(x_{0})-f^{\star}\bigr)\,\varepsilon^{-2},
\qquad
c=\frac{1}{M+G/\Delta_{\min}}.
\]
Thus the dominant constant is
\[
C_{1}= \frac{2\bigl(M+G/\Delta_{\min}\bigr)}{\eta_{1}^{2}}.
\]

\section*{Special Cases}
\begin{enumerate}[leftmargin=*]
\item \textbf{Convex quadratic case.}
If \(f(x)=\frac12 x^{\top}Qx-b^{\top}x\) with \(Q\succeq0\) and we set \(B_{k}=Q\),
the Cauchy point is the exact solution of (TR) for any \(\Delta_{k}\ge
\|Q^{-1}g_{k}\|\).  Hence every iteration is successful, \(\rho_{k}=1\), and the
algorithm converges in a single step to the global minimizer.
\item \textbf{Exact subproblem solution.}
When the subproblem is solved exactly, condition (1) holds with
\(\eta_{1}=1\).  The bound (S7) becomes
\[
f(x_{k})-f(x_{k+1})\ge
\frac{\|g_{k}\|^{2}}{2\bigl(M+\frac{\|g_{k}\|}{\Delta_{k}}\bigr)},
\]
which improves the constant \(C_{1}\) by a factor of \(\eta_{1}^{2}=1\).
\item \textbf{Large‑radius regime.}
If \(\Delta_{k}\ge \|g_{k}\|/M\) for all \(k\), then
\(M+\frac{\|g_{k}\|}{\Delta_{k}}\le 2M\) and the complexity simplifies to
\[
N(\varepsilon)\le
\frac{f(x_{0})-f^{\star}}{\eta_{1}^{2}M}\,\varepsilon^{-2}.
\]
\end{enumerate}

\begin{thebibliography}{9}
\bibitem{conn1999}
A.~R. Conn, N.~I.~M. Gould, and Ph.~L. Toint,
\emph{Trust-Region Methods}, 
SIAM, 1999.
\end{thebibliography}

\end{document}